% GENERATED FILE. Edit canonical lessons, then run npm run generate:books.
% canonical-source-hash: fnv1a64:b07743a21a03cd7c
% canonical-lessons: ML-C184-otiyuka, ML-C184-bhedamakuka, ML-C184-vyayamam-ceyyuka, ML-C184-purippikkuka, ML-C184-oppituka

\chapter{To break, To get better, To exercise, To fill in, To sign}
\label{ch:ml-to-break-to-get-better-to-exercise-to-fill-in-to-sign}
\hlchaptermodality{184}

\begin{chapteropening}
\textbf{By the end of this chapter:} \emph{I can say to break, to get better, to exercise, to fill in and to sign.}

Complete the last lesson of chapter 184: I can say to break, to get better, to exercise, to fill in and to sign.
\end{chapteropening}

\section[\texorpdfstring{\ml{ഒടിയുക} (oṭiyuka)}{ഒടിയുക (oṭiyuka)}]{\ml{ഒടിയുക} (oṭiyuka) — to break (a bone, a stick)}

\label{lesson:ML-C184-otiyuka}



\begin{quote}
Before the new one: say the Malayalam for to get, then the Malayalam for to lose.
\end{quote}

\subsection*{You'll want to know: \ml{ഒടിയുക}}
\textbf{\ml{ഒടിയുക}} — \emph{oṭiyuka} — \textquotedblleft{}to break (a bone, a stick)\textquotedblright{}.

\begin{cousinweb}
\textbf{\ml{ഒടിയുക}} (\emph{oṭiyuka}) is native Dravidian, one word with Tamil \emph{oṭi} and Kannada \emph{oḍi} (DEDR 946).
\end{cousinweb}

\begin{grammarlens}[title={What you've built}]
One more word for this chapter.
\end{grammarlens}

\subsection*{Your turn}
\begin{itemize}
\raggedright
  \item \emph{Say it:} \emph{oṭiyuka}
  \item \emph{Say it:} \emph{oṭiyuka}, once more
  \item \emph{Say it:} say \emph{oṭiyuka}
  \item \emph{Recall:} say the Malayalam for to get, then the Malayalam for to lose, then say \emph{oṭiyuka} again
\end{itemize}

\begin{tcolorbox}[breakable,colback=teal!4,colframe=teal!35!black,title={Before you move on}]
What does \ml{ഒടിയുക} mean? (\textquotedblleft{}To break (a bone, a stick)\textquotedblright{}.) Say it once more.
\end{tcolorbox}
\section[\texorpdfstring{\ml{ഭേദമാകുക} (bhēdamākuka)}{ഭേദമാകുക (bhēdamākuka)}]{\ml{ഭേദമാകുക} (bhēdamākuka) — to get better, to recover}

\label{lesson:ML-C184-bhedamakuka}



\begin{quote}
Before the new one: say the Malayalam for to lose, then the Malayalam for to break.
\end{quote}

\subsection*{You'll want to know: \ml{ഭേദമാകുക}}
\textbf{\ml{ഭേദമാകുക}} — \emph{bhēdamākuka} — \textquotedblleft{}to get better, to recover\textquotedblright{}.

\begin{cousinweb}
\textbf{\ml{ഭേദമാകുക}} (\emph{bhēdamākuka}) joins Sanskrit \emph{bhēda}, \textquotedblleft{}division, difference\textquotedblright{} (CDIAL 9610), to native \emph{ākuka}, \textquotedblleft{}to become\textquotedblright{} (DEDR 333): to get better is \textquotedblleft{}to become different\textquotedblright{}. DEDR 333 compares Sanskrit \emph{ām}.
\end{cousinweb}

\begin{grammarlens}[title={What you've built}]
One more word for this chapter.
\end{grammarlens}

\subsection*{Your turn}
\begin{itemize}
\raggedright
  \item \emph{Say it:} \emph{bhēdamākuka}
  \item \emph{Say it:} \emph{bhēdamākuka}, once more
  \item \emph{Say it:} say \emph{bhēdamākuka}
  \item \emph{Recall:} say the Malayalam for to lose, then the Malayalam for to break, then say \emph{bhēdamākuka} again
\end{itemize}

\begin{tcolorbox}[breakable,colback=teal!4,colframe=teal!35!black,title={Before you move on}]
What does \ml{ഭേദമാകുക} mean? (\textquotedblleft{}To get better, to recover\textquotedblright{}.) Say it once more.
\end{tcolorbox}
\section[\texorpdfstring{\ml{വ്യായാമം} \ml{ചെയ്യുക} (vyāyāmaṁ ceyyuka)}{വ്യായാമം ചെയ്യുക (vyāyāmaṁ ceyyuka)}]{\ml{വ്യായാമം} \ml{ചെയ്യുക} (vyāyāmaṁ ceyyuka) — to exercise}

\label{lesson:ML-C184-vyayamam-ceyyuka}



\begin{quote}
Before the new one: say the Malayalam for to break, then the Malayalam for to get better.
\end{quote}

\subsection*{You'll want to know: \ml{വ്യായാമം} \ml{ചെയ്യുക}}
\textbf{\ml{വ്യായാമം} \ml{ചെയ്യുക}} — \emph{vyāyāmaṁ ceyyuka} — \textquotedblleft{}to exercise\textquotedblright{}.

\begin{cousinweb}
\textbf{\ml{വ്യായാമം} \ml{ചെയ്യുക}} (\emph{vyāyāmaṁ ceyyuka}) joins Sanskrit \emph{vyāyāma}, \textquotedblleft{}exertion\textquotedblright{} (CDIAL 12209), to native \emph{ceyyuka}, \textquotedblleft{}to do\textquotedblright{} (DEDR 1957, with Tamil \emph{cey}).
\end{cousinweb}

\begin{grammarlens}[title={What you've built}]
One more word for this chapter.
\end{grammarlens}

\subsection*{Your turn}
\begin{itemize}
\raggedright
  \item \emph{Say it:} \emph{vyāyāmaṁ ceyyuka}
  \item \emph{Say it:} \emph{vyāyāmaṁ ceyyuka}, once more
  \item \emph{Say it:} say \emph{vyāyāmaṁ ceyyuka}
  \item \emph{Recall:} say the Malayalam for to break, then the Malayalam for to get better, then say \emph{vyāyāmaṁ ceyyuka} again
\end{itemize}

\begin{tcolorbox}[breakable,colback=teal!4,colframe=teal!35!black,title={Before you move on}]
What does \ml{വ്യായാമം} \ml{ചെയ്യുക} mean? (\textquotedblleft{}To exercise\textquotedblright{}.) Say it once more.
\end{tcolorbox}
\section[\texorpdfstring{\ml{പൂരിപ്പിക്കുക} (pūrippikkuka)}{പൂരിപ്പിക്കുക (pūrippikkuka)}]{\ml{പൂരിപ്പിക്കുക} (pūrippikkuka) — to fill in (a form)}

\label{lesson:ML-C184-purippikkuka}



\begin{quote}
Before the new one: say the Malayalam for to get better, then the Malayalam for to exercise.
\end{quote}

\subsection*{You'll want to know: \ml{പൂരിപ്പിക്കുക}}
\textbf{\ml{പൂരിപ്പിക്കുക}} — \emph{pūrippikkuka} — \textquotedblleft{}to fill in (a form)\textquotedblright{}.

\begin{cousinweb}
\textbf{\ml{പൂരിപ്പിക്കുക}} (\emph{pūrippikkuka}) is built on Sanskrit \emph{pūrayati}, \textquotedblleft{}fills\textquotedblright{} (CDIAL 8335), made into a Malayalam causative: \textquotedblleft{}make full\textquotedblright{}.
\end{cousinweb}

\begin{grammarlens}[title={What you've built}]
One more word for this chapter.
\end{grammarlens}

\subsection*{Your turn}
\begin{itemize}
\raggedright
  \item \emph{Say it:} \emph{pūrippikkuka}
  \item \emph{Say it:} \emph{pūrippikkuka}, once more
  \item \emph{Say it:} say \emph{pūrippikkuka}
  \item \emph{Recall:} say the Malayalam for to get better, then the Malayalam for to exercise, then say \emph{pūrippikkuka} again
\end{itemize}

\begin{tcolorbox}[breakable,colback=teal!4,colframe=teal!35!black,title={Before you move on}]
What does \ml{പൂരിപ്പിക്കുക} mean? (\textquotedblleft{}To fill in (a form)\textquotedblright{}.) Say it once more.
\end{tcolorbox}
\section[\texorpdfstring{\ml{ഒപ്പിടുക} (oppiṭuka)}{ഒപ്പിടുക (oppiṭuka)}]{\ml{ഒപ്പിടുക} (oppiṭuka) — to sign}

\label{lesson:ML-C184-oppituka}



\begin{quote}
Before the new one: say the Malayalam for to exercise, then the Malayalam for to fill in.
\end{quote}

\subsection*{You'll want to know: \ml{ഒപ്പിടുക}}
\textbf{\ml{ഒപ്പിടുക}} — \emph{oppiṭuka} — \textquotedblleft{}to sign\textquotedblright{}.

\begin{cousinweb}
\textbf{\ml{ഒപ്പിടുക}} (\emph{oppiṭuka}) is native Dravidian, \textquotedblleft{}put one's signature\textquotedblright{}: \emph{oppu}, \textquotedblleft{}conformity; signature\textquotedblright{} (DEDR 924, with Kannada \emph{oppu}), and \emph{iṭuka}, \textquotedblleft{}to put\textquotedblright{} (DEDR 442, with Tamil \emph{iṭu}).
\end{cousinweb}

\begin{grammarlens}[title={What you've built}]
One more word for this chapter.
\end{grammarlens}

\subsection*{Your turn}
\begin{itemize}
\raggedright
  \item \emph{Say it:} \emph{oppiṭuka}
  \item \emph{Say it:} \emph{oppiṭuka}, once more
  \item \emph{Say it:} say \emph{oppiṭuka}
  \item \emph{Recall:} say the Malayalam for to exercise, then the Malayalam for to fill in, then say \emph{oppiṭuka} again
\end{itemize}

\begin{tcolorbox}[breakable,colback=teal!4,colframe=teal!35!black,title={Before you move on}]
What does \ml{ഒപ്പിടുക} mean? (\textquotedblleft{}To sign\textquotedblright{}.) Say it once more.
\end{tcolorbox}
